\section{Requirements}
\label{sec:requirements}

The requirements were organized with a MoSCoW-style priority structure. Must-have requirements define the minimum useful consultation-support system. Should-have and could-have requirements describe important improvements and possible extensions. Will-not-have requirements clarify the boundaries of the implementation.

Their implementation and validation evidence is traced in Section~\ref{subsec:requirements-coverage}, where each requirement is linked to the report sections and appendices that demonstrate it.

\subsection{Actors and Use Cases}
\label{subsec:actors-use-cases}

This subsection defines the system actors and the required user goals. The design chapter reuses these actors as inputs instead of redefining them, so role responsibilities and use cases remain traceable from one place.

The implemented workflow is centered on three roles: doctors, administrators, and patients. Doctors own the consultation and approval flow, patients can book appointments and read their own consultation and prescription information, and administrators manage configuration that affects AI prompts and email templates. All three roles use the same email/password login screen. Staff and patient passwords are stored as bcrypt hashes, and the current implementation does not include doctor-generated password setup links, magic-login links, or patient login URLs sent by doctors. The doctor email workflow sends approved follow-up material with a PDF attachment rather than an authentication link. Table~\ref{tab:primary-use-cases} summarizes the primary implemented use cases used later as design inputs.

\begin{figure}[htbp]
\centering
\resizebox{\textwidth}{!}{%
\begin{tikzpicture}[
	font=\scriptsize,
	actor/.style={draw=VertexBlue!75, fill=VertexBlue!8, rounded corners=1mm, align=center, minimum width=1.6cm, minimum height=0.58cm},
	usecase/.style={draw=VertexBlue!55, fill=white, ellipse, align=center, text width=2.15cm, minimum height=0.58cm, inner sep=1.4mm},
	authbox/.style={draw=VertexInk, fill=VertexPanel, rounded corners=1.5mm, inner sep=2.2mm},
	doctorbox/.style={draw=VertexInk, fill=VertexTeal!8, rounded corners=1.5mm, inner sep=2.2mm},
	patientbox/.style={draw=VertexInk, fill=VertexBlue!6, rounded corners=1.5mm, inner sep=2.2mm},
	adminbox/.style={draw=VertexInk, fill=VertexRed!5, rounded corners=1.5mm, inner sep=2.2mm},
	service/.style={draw=VertexAmber!75, fill=VertexAmber!10, rounded corners=1mm, align=center, text width=1.65cm, minimum height=0.62cm, inner sep=1.2mm},
	data/.style={draw=VertexTeal!75, fill=VertexTeal!10, cylinder, shape border rotate=90, aspect=0.25, align=center, text width=1.25cm, minimum height=0.68cm, inner sep=1.2mm},
	link/.style={-{Latex[length=1.8mm]}, draw=VertexInk!62, line width=0.45pt},
	assoc/.style={draw=VertexInk!50, line width=0.45pt}
]
\node[font=\bfseries] at (0,5.15) {OPD-VERTEX (MedFlow)};

\node[actor] (doctor) at (-5.35,4.35) {Doctor};
\node[actor] (patient) at (0,4.35) {Patient};
\node[actor] (admin) at (5.35,4.35) {Admin};

\node[usecase, text width=1.25cm] (login) at (-0.80,3.62) {Log In};
\node[usecase, text width=1.25cm] (logout) at (0.80,3.62) {Log Out};
\node[authbox, fit=(login)(logout), label={[font=\bfseries]above:Authentication}] (auth) {};

\node[usecase] (record) at (-6.15,2.92) {Record \& Transcribe Audio};
\node[usecase] (create) at (-3.88,2.92) {Create Consultation};
\node[usecase] (browse) at (-1.62,2.92) {Browse \& Search Patients};
\node[usecase] (viewconsult) at (-6.15,2.18) {View Consultations};
\node[usecase] (review) at (-3.88,2.18) {Review Draft Notes};
\node[usecase] (approve) at (-1.62,2.18) {Approve / Reject Draft};
\node[usecase] (report) at (-6.15,1.44) {Generate Clinical Report};
\node[usecase] (suggest) at (-3.88,1.44) {Suggestive Safety Review};
\node[usecase] (pdf) at (-1.62,1.44) {Generate Prescription PDF};
\node[usecase] (email) at (-5.02,0.70) {Email Prescription};
\node[usecase] (appt) at (-2.75,0.70) {View/Create/Accept Appointments};
\node[doctorbox, fit=(record)(create)(browse)(viewconsult)(review)(approve)(report)(suggest)(pdf)(email)(appt), label={[font=\bfseries]above:Doctor - Clinical Workflow}] (doctorflow) {};

\node[usecase] (profile) at (1.65,2.92) {View Own Profile};
\node[usecase] (ownrx) at (3.92,2.92) {View Own Prescriptions};
\node[usecase] (ownappt) at (2.78,2.18) {View/Create/Delete Appointments};
\node[patientbox, fit=(profile)(ownrx)(ownappt), label={[font=\bfseries]above:Patient Features}] (patientfeatures) {};

\node[usecase] (prompts) at (4.35,1.42) {Manage Prompt Templates};
\node[usecase] (templates) at (6.62,1.42) {Manage Email Templates};
\node[usecase] (health) at (5.48,0.68) {View System Health};
\node[adminbox, fit=(prompts)(templates)(health), label={[font=\bfseries]above:Admin}] (adminfeatures) {};

\node[service] (ollama) at (-6.15,-0.52) {Ollama LLM\\Report Gen};
\node[service] (whisper) at (-3.78,-0.52) {Whisper API\\Speech-to-Text};
\node[service] (reportlab) at (-1.42,-0.52) {ReportLab\\PDF};
\node[service] (smtp) at (0.92,-0.52) {SMTP\\Email};
\node[data] (mysql) at (3.35,-0.52) {MySQL};
\node[data] (mongo) at (5.45,-0.52) {MongoDB};

\draw[assoc] (doctor) -- (auth);
\draw[assoc] (patient) -- (auth);
\draw[assoc] (admin) -- (auth);
\draw[assoc] (doctor) -- (doctorflow.north);
\draw[assoc] (patient) -- (patientfeatures.north);
\draw[assoc] (admin) -- (adminfeatures.north);
\draw[link] (record) -- (whisper);
\draw[link] (report) -- (ollama);
\draw[link] (suggest) -- (ollama);
\draw[link] (pdf) -- (reportlab);
\draw[link] (email) -- (smtp);
\draw[link] (appt) -- (mysql);
\draw[link] (ownappt) -- (mysql);
\draw[link] (prompts) -- (mongo);
\draw[link] (templates) -- (mongo);

\node[anchor=north west, align=left, font=\tiny] at (-7.32,-1.08) {\textbf{Legend:} actors = role boxes; ovals = use cases; rounded rectangles = functional modules; amber boxes = external/service adapters; cylinders = databases; arrows = direct service or data dependency.};
\end{tikzpicture}%
}
\caption{Implemented Med Flow use-case view for doctors, patients, administrators, infrastructure components, and external services.}
\label{fig:use-case-diagram}
\end{figure}
\FloatBarrier

\begin{table}[htbp]
\centering
\small
\begin{tabularx}{\textwidth}{@{}P{0.10\textwidth}P{0.16\textwidth}YP{0.30\textwidth}@{}}
\TableTopRule
\rowcolor{VertexTableBlue}
\VertexTableHead{ID} & \VertexTableHead{Actor} & \VertexTableHead{Goal} & \VertexTableHead{Primary flow} \\
\TableMidRule
\rowcolor{VertexTableStripe}
UC-1 & All roles & Authenticate and reach the correct dashboard. & User enters email and password, receives a JWT cookie, and is routed to the role-aware dashboard. \\
UC-2 & Patient & Book an appointment. & Patient selects appointment details for themselves, submits reason and time, and creates appointment context. \\
\rowcolor{VertexTableStripe}
UC-3 & Doctor / Admin & Manage appointment status. & Doctor or administrator confirms, cancels, or marks no-show appointments; only doctors start consultations. \\
UC-4 & Doctor & Record and transcribe a consultation. & Doctor starts the linked consultation, records audio, receives transcript updates, and finalizes transcription. \\
\rowcolor{VertexTableStripe}
UC-5 & Doctor & Generate AI-supported clinical material. & Doctor sends transcript and patient context through local AI services to create report, prescription, and alert drafts. \\
UC-6 & Doctor & Review and edit generated drafts. & Doctor reads transcript, structured report, prescription draft, and suggestive alerts, then edits Markdown if needed. \\
\rowcolor{VertexTableStripe}
UC-7 & Doctor & Approve, export, and send final records. & Doctor approves or rejects drafts, downloads the PDF, and can send approved follow-up email with attachment. \\
UC-8 & Doctor & Search patient history. & Doctor searches patients, opens a patient detail page, and reviews consultation and prescription history. \\
\rowcolor{VertexTableStripe}
UC-9 & Patient & View own clinical information. & Patient logs in with email/password and sees only their own appointments, consultations, and prescriptions. \\
UC-10 & Administrator & Maintain configurable content and auditability. & Administrator edits prompts and email templates and reviews audit entries for those changes. \\
\TableBottomRule
\end{tabularx}
\caption{Primary Med Flow use cases by actor.}
\label{tab:primary-use-cases}
\end{table}
\FloatBarrier

\subsection{Must-Have Requirements}
\label{subsec:must-have-requirements}

\begingroup
\small
\begin{longtable}{@{}P{0.16\textwidth}P{0.76\textwidth}@{}}
\caption{Must-have requirements.}\\
\TableTopRule
\rowcolor{VertexTableBlue}
\VertexTableHead{ID} & \VertexTableHead{Requirement} \\
\TableMidRule
\rowcolor{VertexTableStripe}
F-1 & The system must manage a doctor consultation queue, including consultation context created from booking or scheduling. \\
F-2 & The system must capture consultation audio locally. \\
\rowcolor{VertexTableStripe}
F-3 & The system must transcribe speech to text locally. \\
F-4 & The system must generate a doctor-editable structured clinical note from the transcript and patient context, including presenting problem, relevant history, assessment, treatment or medication plan, follow-up, and patient instructions when supported by the consultation. \\
\rowcolor{VertexTableStripe}
F-5 & The system must generate a doctor-editable prescription draft with medication name, dosage, frequency, duration, patient instructions, and review alerts for missing dosage information or possible allergy/contraindication conflicts. \\
F-6 & The system must allow doctors to edit, approve, or reject generated clinical notes and prescription drafts before they become final records. \\
\rowcolor{VertexTableStripe}
F-7 & The system must generate downloadable PDFs from approved clinical report and prescription data. \\
F-8 & The system must send approved patient follow-up or prescription summaries by configured email and prevent unapproved drafts from being emailed. \\
\rowcolor{VertexTableStripe}
NF-1 & Privacy: all AI processing must be performed locally. \\
NF-2 & Security: the system must use authentication and role-based access for doctors, administrators, and patients, with patients restricted to their own records. \\
\rowcolor{VertexTableStripe}
NF-3 & Data protection: the system must hash stored passwords and be designed for encrypted deployment channels and storage. \\
NF-4 & Modularity: components must be loosely coupled and replaceable. \\
\rowcolor{VertexTableStripe}
NF-5 & Reliability: the system must persist consultation state, transcripts, generated documents, and prescription versions so interrupted workflows can be retried without losing approved data. \\
NF-6 & Auditability: the system must preserve traceability for important actions and edits, including prompt updates, email-template updates, doctor report edits with version history, and review approval or rejection status changes. \\
\TableBottomRule
\end{longtable}
\endgroup

\subsection{Should-Have Requirements}
\label{subsec:should-have-requirements}

\begingroup
\small
\begin{longtable}{@{}P{0.16\textwidth}P{0.76\textwidth}@{}}
\caption{Should-have requirements.}\\
\TableTopRule
\rowcolor{VertexTableBlue}
\VertexTableHead{ID} & \VertexTableHead{Requirement} \\
\TableMidRule
\rowcolor{VertexTableStripe}
F-9 & The system should display a live transcript during consultation. \\
F-10 & The system should allow doctors to view past consultations. \\
\rowcolor{VertexTableStripe}
F-12 & The system should keep booking and scheduling replaceable, so an external clinic or public-sector booking service can create consultation context through stable interfaces. \\
NF-7 & Usability: the interface should be user-friendly and efficient. \\
NF-8 & Performance: on the target machine, local transcription should operate with RTF $\leq 1.0$. For a typical consultation transcript, the system should generate structured notes, prescription draft, and clinical alerts within $\leq 40$ seconds on recommended hardware and within $\leq 120$ seconds on minimum hardware. \\
\TableBottomRule
\end{longtable}
\endgroup

\subsection{Could-Have Requirements}
\label{subsec:could-have-requirements}

\begingroup
\small
\begin{longtable}{@{}P{0.16\textwidth}P{0.76\textwidth}@{}}
\caption{Could-have requirements.}\\
\TableTopRule
\rowcolor{VertexTableBlue}
\VertexTableHead{ID} & \VertexTableHead{Requirement} \\
\TableMidRule
\rowcolor{VertexTableStripe}
F-11 & The system could support multiple doctors for sharing documents, such as a shareable report sent to a doctor other than the family doctor. \\
NF-9 & The system could allow administrators to configure how long consultation data, transcripts, and prescriptions are stored, and must support secure deletion of expired records. \\
\TableBottomRule
\end{longtable}
\endgroup

\subsection{Will-Not-Have Requirements}
\label{subsec:will-not-have-requirements}

\begingroup
\small
\begin{longtable}{@{}P{0.16\textwidth}P{0.76\textwidth}@{}}
\caption{Will-not-have requirements.}\\
\TableTopRule
\rowcolor{VertexTableBlue}
\VertexTableHead{ID} & \VertexTableHead{Requirement} \\
\TableMidRule
\rowcolor{VertexTableStripe}
NF-10 & The system will not use cloud-based AI services. \\
NF-11 & The system alone, excluding the doctor, is not enough for diagnosis. \\
\TableBottomRule
\end{longtable}
\endgroup
